% problem-000089 5 分光光度分析
\section*{5. 分光光度分析}
（图：）

氧化得到的偶联产物在550 nm波⻓处的摩尔消光系数为 $1 . 8 \times 1 0 ^ { 4 } \mathrm { L m o l ^ { - 1 } c m ^ { - 1 } }$ 。利⽤这⼀反应，通过分光光度法可测定起始溶液中的银离⼦浓度。

${ \mathrm { \bf 5 - 1 } } \quad 2 5 \ { \mathrm { \bf ~ ^ { \circ } C } } ,$ ，在纯⽔中配得⼄酸银的饱和⽔溶液。取1.00 mL溶液，按计量加⼊两种苯⼆胺，⾄显⾊稳定后，稀释定容⾄ $1 \mathrm { L } _ { \circ }$ 。⽤1 cm的⽐⾊⽫，测得溶液在550 nm处的吸光度为0.225。计算⼄酸银的溶度积 $\updownarrow K _ { \mathsf { s p } ^ { \mathsf { c } } }$ 。（提⽰：可合理忽略副反应。朗伯-⽐尔⽅程 $A = \varepsilon l c$ ，�吸光度；�摩尔消光系数；�光通过溶液的路径⻓；�溶液浓度。）

\subsection*{5-2}
室温下将⼄酸银与4-三甲铵基苯硫醇的六氟磷酸盐(HL+ $) ( \mathrm { P F } _ { 6 } ^ { - } ) \{ \mathrm { H L } ^ { + } = [ \mathrm { H S C } _ { 6 } \mathrm { H } _ { 4 } \mathrm { N } ( \mathrm { C H } _ { 3 } ) _ { 3 } ] ^ { + }$ +}混合研磨，两种固体发⽣反应，得到⼀种淡⻩⾊粉末。粉末⽤⼄腈处理，结晶出⽆⾊溶剂合物A，化学式为

$\mathrm { \Delta [ A g _ { 9 } L _ { 8 } ( C H _ { 3 } C N ) _ { \it x } ] ( P F _ { 6 } ) _ { \it y } ] { \cdot } \it { n C H _ { 3 } C N _ { \circ } } }$ 元素分析结果给出，A中C, H, N, P的含量（质量分数，%）分别为：C27.40，H 3.35，N 6.24，P 6.90。室温下在真空中⼩⼼处理A，A可脱去外界的溶剂分⼦，失重 $2 . 0 \% _ { \mathsf { c } }$

\subsection*{5-2}
-1 通过计算，确定A的化学式。（提⽰：元素分析中，氢含量测定误差常常较⼤）

\subsection*{5-2}
-2 写出两种固体发⽣反应⽣成淡⻩⾊粉末的⽅程式（要求系数为最简整数⽐）。
